% =============================================================================
% SOM_dune_development_appendix.tex
% Appendix: Supplementary Figures and Tables
% =============================================================================

\clearpage
\appendix
\section{Supplementary Figures}
\label{sec:appendix_figures}

% --- Detailed Substrate Characterisation ---
\begin{figure}[htbp]
    \centering
    \begin{subfigure}[b]{0.85\textwidth}
        \centering
        \includegraphics[width=\textwidth]{substrate_tn.pdf}
        \caption{Total nitrogen by substrate}
        \label{fig:app_substrate_tn}
    \end{subfigure}
    \par\vspace{1em}
    \begin{subfigure}[b]{0.85\textwidth}
        \centering
        \includegraphics[width=\textwidth]{substrate_in_out_comparison.pdf}
        \caption{Inside vs.\ outside sampling comparison}
        \label{fig:app_substrate_inout}
    \end{subfigure}
    \caption{Additional substrate characterisation: (a) Total nitrogen content across the 10 substrates, confirming the nutrient gradient, and (b) comparison of physical and chemical properties between inside- and outside-box sampling positions within the alternative substrates, validating spatial homogeneity ($p > 0.7$ for all Wilcoxon rank-sum tests; \cref{tab:substrate_stats}).}
    \label{fig:app_substrate_variance}
\end{figure}

% --- Detailed Aboveground Distributions ---
\begin{figure}[htbp]
    \centering
    \begin{subfigure}[b]{0.85\textwidth}
        \centering
        \includegraphics[width=\textwidth]{marram_height_all.pdf}
        \caption{Marram height distributions by time point}
        \label{fig:app_height_all}
    \end{subfigure}
    \par\vspace{1em}
    \begin{subfigure}[b]{0.85\textwidth}
        \centering
        \includegraphics[width=\textwidth]{number_of_leaves.pdf}
        \caption{Leaf count distributions by time point}
        \label{fig:app_leaves_all}
    \end{subfigure}
    \caption{Detailed time-series boxplots of aboveground plant performance by substrate and measurement date. These non-parametric distributions underpin the model-predicted trajectories (Gaussian LMM for height, Poisson GLMM for leaf count) shown in the main text (\cref{fig:aboveground_growth}).}
    \label{fig:app_aboveground_dist}
\end{figure}

% --- Belowground Biomass ---
\begin{figure}[htbp]
    \centering
    \begin{subfigure}[b]{0.85\textwidth}
        \centering
        \includegraphics[width=\textwidth]{belowground_main.pdf}
        \caption{Root biomass and root:shoot ratio}
        \label{fig:app_root_shoot}
    \end{subfigure}
    \par\vspace{1em}
    \begin{subfigure}[b]{0.85\textwidth}
        \centering
        \includegraphics[width=\textwidth]{height_vs_shoot_biomass.pdf}
        \caption{Height vs.\ shoot biomass ($r = 0.45$, $p = 0.001$)}
        \label{fig:app_height_biomass}
    \end{subfigure}
    \par\vspace{1em}
    \begin{subfigure}[b]{0.85\textwidth}
        \centering
        \includegraphics[width=\textwidth]{root_depth_profile.pdf}
        \caption{Root biomass depth profile}
        \label{fig:app_root_depth}
    \end{subfigure}
    \par\vspace{1em}
    \begin{subfigure}[b]{0.85\textwidth}
        \centering
        \includegraphics[width=\textwidth]{root_vs_shoot_scatter.pdf}
        \caption{Root vs.\ shoot biomass by substrate}
        \label{fig:app_root_vs_shoot}
    \end{subfigure}
    \caption{Belowground biomass allocation. No significant substrate effects were detected by either Gamma GLMs or Kruskal--Wallis tests for total root biomass ($F = 0.68$, $p = 0.72$; $\chi^2 = 6.87$, $p = 0.65$), root:shoot ratio ($F = 1.41$, $p = 0.22$; $\chi^2 = 9.67$, $p = 0.38$), or shoot biomass ($F = 0.43$, $p = 0.91$; $\chi^2 = 5.76$, $p = 0.76$). Height was significantly positively correlated with shoot biomass (b), confirming the allometric relationship used in the Living Dunes vegetation model.}
    \label{fig:app_belowground}
\end{figure}

% --- Detailed Microclimate Data ---
\begin{figure}[htbp]
    \centering
    \includegraphics[width=0.85\textwidth]{microclimate_boxplots.pdf}
    \caption{Overview of microclimate variables recorded by TMS-4 loggers during the germination experiment ($\sim$225 days). Soil temperature did not differ significantly among substrates ($\chi^2(5) = 1.23$, $p = 0.94$; Kruskal--Wallis), while volumetric soil moisture showed a clear SOM gradient ($\chi^2(5) = 634.92$, $p < 2 \times 10^{-16}$; Kruskal--Wallis).}
    \label{fig:app_clim_vars}
\end{figure}

\begin{figure}[htbp]
    \centering
    \begin{subfigure}[b]{0.85\textwidth}
        \centering
        \includegraphics[width=\textwidth]{soil_temperature_timeseries.pdf}
        \caption{Soil temperature time series}
        \label{fig:app_soil_temp}
    \end{subfigure}
    \par\vspace{1em}
    \begin{subfigure}[b]{0.85\textwidth}
        \centering
        \includegraphics[width=\textwidth]{diurnal_temperature.pdf}
        \caption{Diurnal temperature range by substrate}
        \label{fig:app_temp_range}
    \end{subfigure}
    \caption{Temperature dynamics during the germination experiment: (a) daily mean soil temperature showing seasonal warming from April to August, with no significant substrate differences, and (b) diurnal temperature amplitude highlighting minor substrate-specific thermal buffering.}
    \label{fig:app_temperature_logs}
\end{figure}

\begin{figure}[htbp]
    \centering
    \includegraphics[width=0.65\textwidth]{soil_moisture_timeseries.pdf}
    \caption{Volumetric soil moisture time series by substrate, demonstrating the persistent moisture gradient attributable to the SOM ``sponge effect''. The dark substrate (D) maintained consistently higher moisture ($\overline{m} = 19.2\%$) than the reference beach sand (R, $\overline{m} = 10.8\%$).}
    \label{fig:app_moisture_timeseries}
\end{figure}


% --- Supplementary Tables ---
\clearpage
\section{Supplementary Tables}
\label{sec:appendix_tables}

% --- Substrate characterisation tables ---
\begin{table}[htbp]
\centering
\caption{Wilcoxon rank-sum test results comparing sediment properties between alternative (dredged) and control (beach) substrates, and between inside and outside sampling positions within the alternative substrate.}
\label{tab:substrate_stats}
\small
\begin{tabular}{l l r r}
\toprule
Comparison & Property & $W$ & $p$-value \\
\midrule
Alternative vs Control & TOC (\%) & 55 & 1.64e-08 \\
Alternative vs Control & TN (\%) & 176 & 1.34e-05 \\
Alternative vs Control & $D_{50}$ ($\mu$m) & 840 & 0.00634 \\
\midrule
Inside vs Outside & TOC (\%) & 517 & 0.75 \\
Inside vs Outside & TN (\%) & 562 & 0.808 \\
Inside vs Outside & $D_{50}$ ($\mu$m) & 560 & 0.828 \\
\bottomrule
\end{tabular}
\end{table}


\begin{table}[htbp]
\centering
\caption{Physical and chemical properties of the 10 experimental substrates used in the common garden experiment.}
\label{tab:substrate_summary}
\small
\begin{tabular}{l r r r l}
\toprule
Substrate & TOC (\%) & TN (\%) & $D_{50}$ ($\mu$m) & Description \\
\midrule
Ref          & 0.06 & 0.036 & 267 & Control (Natural Sand) \\
White        & 0.25 & 0.073 & 188 & Pure White Dredged Material \\
12.5% D      & 0.49 & 0.087 & 187 & Mixture \\
25% D        & 0.73 & 0.101 & 187 & Mixture \\
37.5% D      & 0.97 & 0.116 & 186 & Mixture \\
50% D        & 1.21 & 0.130 & 185 & Highest Growth Potential \\
62.5% D      & 1.45 & 0.144 & 185 & Mixture \\
75% D        & 1.68 & 0.159 & 184 & Mixture \\
87.5% D      & 1.92 & 0.173 & 184 & Mixture \\
Dark         & 2.16 & 0.187 & 183 & Pure Dark Dredged Material \\
\bottomrule
\end{tabular}
\end{table}


% --- Aboveground growth tables ---
\begin{table}[htbp]
\centering
\caption{Estimated marginal means (EMMs) for aboveground performance metrics. Height was modelled with a Gaussian LMM (quadratic time); leaf count with a Poisson GLMM (\texttt{glmmTMB}, quadratic time); both with plant identity as random intercept. Groups sharing the same letter are not significantly different (FDR-adjusted pairwise contrasts, $p < 0.05$).}
\label{tab:aboveground_lmm}
\small
\begin{tabular}{l rr cc}
\toprule
Substrate & Height EMM (cm) & Leaf count EMM & Height group & Leaf group \\
\midrule
R        & 43.9 & 8.8 & a & ab \\
D12.5    & 49.8 & 6.7 & ab & a \\
W        & 51.1 & 7.5 & ab & ab \\
D62.5    & 51.4 & 8.8 & ab & ab \\
D        & 53.3 & 12.2 & ab & b \\
D75      & 55.4 & 10.5 & b & ab \\
D87.5    & 56.6 & 10.1 & b & ab \\
D25      & 56.8 & 8.4 & b & ab \\
D37.5    & 57.7 & 8.1 & b & ab \\
D50      & 59.4 & 8.9 & b & ab \\
\bottomrule
\end{tabular}
\end{table}


% --- Belowground biomass tables ---
\begin{table}[htbp]
\centering
\caption{Statistical tests for belowground biomass metrics across the SOM gradient. Gamma GLM (log link) was fitted where data were strictly positive; Kruskal-Wallis is reported in parallel as a non-parametric reference. Benjamini-Hochberg correction applied to pairwise comparisons.}
\label{tab:belowground_stats}
\small
\begin{tabular}{l l r r r}
\toprule
Response & Method & Statistic & df & $p$-value \\
\midrule
Total root biomass & KW & $\chi^2$ = 6.87 & 9 & 0.65 \\
                   & Gamma & F = 0.68 & 9 & 0.724 \\
Root:shoot ratio & KW & $\chi^2$ = 9.67 & 9 & 0.377 \\
                 & Gamma & F = 1.41 & 9 & 0.219 \\
Shoot biomass & KW & $\chi^2$ = 5.76 & 9 & 0.764 \\
              & Gamma & F = 0.43 & 9 & 0.913 \\
\bottomrule
\end{tabular}
\end{table}


\begin{table}[htbp]
\centering
\caption{Belowground biomass and root:shoot ratio per substrate (mean $\pm$ SD). Statistical tests: Gamma GLM (log link) where feasible; Kruskal-Wallis rank-sum otherwise. See \cref{tab:belowground_stats} for test details.}
\label{tab:belowground_summary}
\small
\begin{tabular}{l r rr rr rr}
\toprule
Substrate & $n$ & Root (g) & & Shoot (g) & & R:S ratio & \\
          &     & Mean & SD & Mean & SD & Mean & SD \\
\midrule
R        & 2 & 4.49 & 0.01 & 4.12 & 0.35 & 1.09 & 0.09 \\
W        & 3 & 6.76 & 3.64 & 4.74 & 1.64 & 1.37 & 0.30 \\
D12.5    & 6 & 7.51 & 5.63 & 4.79 & 2.32 & 1.51 & 0.62 \\
D25      & 5 & 5.96 & 0.92 & 5.42 & 1.02 & 1.12 & 0.22 \\
D37.5    & 6 & 7.12 & 3.10 & 5.89 & 1.56 & 1.23 & 0.52 \\
D50      & 6 & 6.05 & 2.46 & 5.79 & 2.32 & 1.12 & 0.39 \\
D62.5    & 6 & 5.98 & 0.80 & 4.96 & 1.45 & 1.29 & 0.37 \\
D75      & 6 & 5.83 & 3.03 & 4.83 & 2.59 & 1.42 & 0.60 \\
D87.5    & 5 & 4.43 & 2.44 & 5.41 & 2.54 & 0.78 & 0.49 \\
D        & 2 & 4.37 & 0.59 & 5.87 & 0.55 & 0.75 & 0.17 \\
\bottomrule
\end{tabular}
\end{table}


% --- Germination tables ---
\begin{table}[htbp]
\centering
\caption{Germination and seedling establishment across substrates. Germination at $t_1$ (10 seeds per pot, $n=16$) modelled with a binomial GLM (quasibinomial due to overdispersion); seedling count over time modelled with a Negative Binomial GLMM (\texttt{glmmTMB}, quadratic time, random intercept per pot). Groups sharing the same letter are not significantly different (FDR-adjusted pairwise contrasts, $p < 0.05$).}
\label{tab:germination_stats}
\small
\begin{tabular}{l r rr rr cc}
\toprule
Substrate & $n$ & Germ. \% & SD & EMM count & SE & Germ. group & Count group \\
\midrule
R      & 16 & 6.2 & 13.6 & 0.4 & 0.1 & a & a \\
W      & 16 & 53.1 & 21.8 & 3.4 & 0.7 & c & bc \\
D25    & 16 & 46.9 & 18.2 & 4.0 & 0.7 & c & bc \\
D50    & 16 & 49.4 & 22.6 & 3.7 & 0.7 & c & bc \\
D75    & 16 & 54.4 & 25.3 & 4.7 & 0.8 & c & c \\
D      & 16 & 23.1 & 22.4 & 2.2 & 0.5 & b & b \\
\bottomrule
\end{tabular}
\end{table}


% --- Microclimate tables ---
\begin{table}[htbp]
\centering
\caption{Microclimate conditions in germination experiment substrates, measured by TMS-4 loggers. Values are daily means $\pm$ SD.}
\label{tab:microclimate}
\small
\begin{tabular}{l r l l l}
\toprule
Substrate & Days & Soil temp. ($^{\circ}$C) & Air temp. ($^{\circ}$C) & Vol. moisture (\%) \\
\midrule
R      & 224 & 18.6 +/- 4.5 & 17.2 +/- 4.0 & 10.8 +/- 3.9 \\
W      & 225 & 18.7 +/- 4.5 & 17.1 +/- 4.1 & 9.1 +/- 3.8 \\
D25    & 221 & 18.8 +/- 4.4 & 17.3 +/- 4.0 & 13.3 +/- 4.8 \\
D50    & 225 & 18.5 +/- 4.6 & 17.2 +/- 4.1 & 14.9 +/- 5.7 \\
D75    & 225 & 18.4 +/- 4.4 & 17.3 +/- 4.0 & 18.1 +/- 7.0 \\
D      & 225 & 18.6 +/- 4.6 & 17.2 +/- 4.0 & 19.2 +/- 7.2 \\
\bottomrule
\end{tabular}
\end{table}


% --- SOM response model tables ---
\begin{table}[htbp]
\centering
\caption{Non-linear regression coefficients for the Gaussian growth response models (Levenberg--Marquardt algorithm). The response function is $y = b_0 + a \cdot \exp\left(-\frac{(\mathrm{SOM} - \mu)^2}{2\sigma^2}\right)$, where $b_0$ is the baseline multiplier, $a$ is the amplitude, $\mu$ is the optimal SOM (\% TOC), and $\sigma$ is the standard deviation of the response. Response variables are normalised to the reference beach sand (R). $\dagger$~Including outlier plant 6b\_5 (97 leaves). $\ddagger$~Outlier removed; amplitude non-significant.}
\label{tab:growth_coefficients}
\small
\begin{tabular}{ll rrrr rr r}
\toprule
Response & Parameter & Estimate & SE & $t$ & $p$ & $R^2$ & $R^2_{adj}$ & $n$ \\
\midrule
Height & $b0$ & 0.5000 & 4.3047 & 0.12 & 0.911 & 0.659 & 0.387 & 120 \\
 & $amp$ & 0.8070 & 4.2818 & 0.19 & 0.857 &  &  &  \\
 & $mu$ & 1.3227 & 0.1441 & 9.18 & $<$0.001 &  &  &  \\
 & $sigma$ & 1.5181 & 4.7023 & 0.32 & 0.758 &  &  &  \\
\midrule
Density$^\dagger$ & $b0$ & 1.0331 & 0.0348 & 29.65 & $<$0.001 & 0.639 & 0.349 & 120 \\
 & $amp$ & 1.4019 & 1.6486 & 0.85 & 0.428 &  &  &  \\
 & $mu$ & 0.6480 & 0.0689 & 9.41 & $<$0.001 &  &  &  \\
 & $sigma$ & 0.1000 & 0.0893 & 1.12 & 0.305 &  &  &  \\
\midrule
Density (clean)$^\ddagger$ & $b0$ & 1.0447 & 0.0357 & 29.24 & $<$0.001 & 0.532 & 0.157 & 119 \\
 & $amp$ & 0.2330 & 0.1374 & 1.70 & 0.141 &  &  &  \\
 & $mu$ & 0.5014 & 0.2245 & 2.23 & 0.067 &  &  &  \\
 & $sigma$ & 0.1000 & 0.1082 & 0.92 & 0.391 &  &  &  \\
\midrule
Root & $b0$ & 0.8858 & 0.8063 & 1.10 & 0.314 & 0.544 & 0.18 & 228 \\
 & $amp$ & 0.6523 & 0.7566 & 0.86 & 0.422 &  &  &  \\
 & $mu$ & 1.2423 & 0.1555 & 7.99 & $<$0.001 &  &  &  \\
 & $sigma$ & 0.6773 & 0.7599 & 0.89 & 0.407 &  &  &  \\
\bottomrule
\end{tabular}
\end{table}


\begin{table}[htbp]
\centering
\caption{Binomial GLMM coefficients for the survival--SOM relationship. The model is $\mathrm{logit}(P_{\mathrm{survival}}) = \beta_0 + \beta_1 \cdot \mathrm{SOM} + u_j$, where $u_j \sim \mathcal{N}(0, \sigma^2_u)$ is a random intercept for spatial block $j$ (column position in common garden). McFadden pseudo-$R^2$:  0 .}
\label{tab:mortality_coefficients}
\begin{tabular}{lrrrr}
\toprule
Component & Estimate & SE & $z$ & $p$ \\
\midrule
\multicolumn{5}{l}{\textit{Fixed effects}} \\
$\beta_0$ (Intercept) & 1.3236 & 0.4252 & 3.11 & 0.0019 \\
$\beta_1$ (SOM) & 0.0105 & 0.3314 & 0.03 & 0.9748 \\
\midrule
\multicolumn{5}{l}{\textit{Random effects}} \\
$\sigma_u$ (column) & 0.0000 & & & \\
\midrule
\multicolumn{5}{l}{\textit{Derived Living Dunes parameters}} \\
Steepness ($-\beta_1$) & -0.0105 & & & \\
Threshold ($-\beta_0/\beta_1$) & -126.21 & & & \\
\bottomrule
\end{tabular}
\end{table}


\begin{table}[htbp]
\centering
\caption{Model fit statistics for the SOM response curve derivations. Growth models use non-linear least squares with the Levenberg--Marquardt algorithm (\texttt{nlsLM}). The mortality model uses a binomial generalised linear mixed model (GLMM) with a random intercept for spatial block.}
\label{tab:model_fit_statistics}
\begin{tabular}{l r r r r r}
\toprule
Model & Log-Lik & AIC & $\sigma_\varepsilon$ & $R^2$ & $n$ \\
\midrule
Height NLS & 14.6 & -19.3 & 0.4022 & 0.659 & 120 \\
Density NLS & 10.5 & -10.9 & 0.5916 & 0.639 & 120 \\
Root NLS & 4.0 & 2.0 & 0.5492 & 0.544 & 228 \\
Mortality GLMM & -61.4 & 128.8 & --- & 0.000$^\dagger$ & 120 \\
\bottomrule
\multicolumn{6}{l}{\footnotesize $\dagger$ McFadden pseudo-$R^2$ for the binomial GLMM.} \\
\end{tabular}
\end{table}

